\section{Interaction Contract}
\label{sec:contract}

\subsection{Failure model}

The contract targets failures that can remain hidden behind a successful
language-model response.  Table~\ref{tab:failure-model} separates these failures
by the relation that becomes inconsistent.  The categories are not claims about
their prevalence; they define the boundary implemented and tested in this work.

\begin{table*}[t]
  \caption{Cross-layer failures addressed by the executable contract.}
  \label{tab:failure-model}
  \small
  \begin{tabularx}{\textwidth}{@{}p{22mm}Y Y Y@{}}
    \toprule
    Failure & Example & Contract decision & Observable consequence \\
    \midrule
    Target identity
      & A collider reports an unknown entity, an obsolete model hash, or a
        source identifier shared by several objects.
      & Reject before routing.
      & No response-model call; target relation and stage are reported. \\
    Target ambiguity
      & The client reports several candidates or a non-resolved deictic state.
      & Reject rather than selecting a candidate.
      & The interface remains in an unresolved state. \\
    Responsibility drift
      & The active agent is not the unique owner and has no declared direct
        handoff to the owner.
      & Resolve asset, then group, then zone ownership; require one owner at the
        highest matching tier and at most one direct handoff.
      & Local/direct routes proceed; transitive or unreachable routes stop. \\
    Action shortcut
      & A scenario or client attempts to invoke an endpoint without a matching
        provider-specific capability and runtime binding.
      & Require an exact scenario-step--capability-use--capability--binding--action
        path.
      & Zero or multiple exact mappings fail closed. \\
    Evidence mismatch
      & A response names a different target, provider, binding, or action from
        the pending request.
      & Reject presentation and retain the failed relation.
      & A transport-level success is not treated as contract success. \\
    Undeclared handoff
      & The response model proposes an agent that is not an allowed handoff
        target.
      & Validate the proposed response before commit.
      & Provisional history is rolled back. \\
    \bottomrule
  \end{tabularx}
\end{table*}

\subsection{Contract objects and executable path}

Let a contract-grounded interaction be represented by the tuple
\[
\mathcal{I} =
\langle m,s,t,p,u,c,b,a,q,o\rangle ,
\]
where \(m\) is a pinned model version and content hash, \(s\) is a scenario
step, \(t\) is the spatial target, \(p\) is the provider, \(u\) is the
capability use in that scenario, \(c\) is the reusable capability, \(b\) is a
runtime binding, \(a\) is a concrete runtime action, \(q\) is the set of
resulting assertions, and \(o\) is runtime observation evidence.  The required
declaration path is
\[
s \rightarrow u \rightarrow c \rightarrow b \rightarrow a.
\]
The scenario step cannot call \(a\) directly.  The capability use names both
\(p\) and its targets, and the executable profile requires exactly one
capability and provider for the evaluated interaction.  Assertions are linked
to a validation case that covers the same use case and executable binding.

Spatial responsibility is resolved independently of display names.  For target
\(t\), the resolver collects candidate agents at three tiers:
\[
R(t) =
\begin{cases}
R_{\mathrm{asset}}(t), & \text{if non-empty},\\
R_{\mathrm{group}}(t), & \text{otherwise if non-empty},\\
R_{\mathrm{zone}}(t), & \text{otherwise}.
\end{cases}
\]
Exactly one agent at the highest active tier is required.  More than one is
ambiguous; none is unassigned.  If the responsible provider is not the active
agent, the current online contract permits exactly one declared handoff edge.
Longer graph paths remain useful offline for diagnosing reachability but are
not traversed within one request.

\subsection{Admission and response conditions}

A deictic request is admitted only if all of the following hold:

\begin{enumerate}
  \item the session is pinned to \(m\), and the supplied model hash equals the
  current project hash and stored contract fingerprint;
  \item the supplied entity and source identifier resolve to the same unique
  scene asset in \(m\), with a resolved, non-ambiguous selection state;
  \item the responsibility rule yields one provider \(p\), and \(p\) is either
  the active agent or one permitted direct handoff away;
  \item one capability use \(u\) for the trusted target and provider resolves to
  one capability \(c\), binding \(b\), and action \(a\);
  \item the request conforms to the declared action schema.
\end{enumerate}

Conditions one through five are request-decidable and are checked before a
response-model call and retained state mutation.  The generated proposal is
then checked for response-dependent conditions: any proposed handoff must be
declared, and the returned model, target, provider, capability, binding, and
action identifiers must agree with the preselected path.  A violation after
generation causes rollback of the provisional conversation state.  Later
runtime events can evaluate assertions as pass, fail, inconclusive, or error,
depending on available observations.

Figure~\ref{fig:admission-boundary} shows the implemented order.  The sequence
matters: placing a schema check after generation would still allow a stale
selection to consume a model call, while trusting a client-supplied provider
would let the caller bypass responsibility resolution.

\begin{figure*}[t]
  \centering
  \begin{tikzpicture}[
    font=\footnotesize,
    stage/.style={
      draw,
      rounded corners=2pt,
      align=center,
      text width=28mm,
      minimum height=17mm,
      inner sep=2mm
    },
    gate/.style={
      stage,
      double,
      double distance=0.6pt
    },
    reject/.style={
      draw,
      rounded corners=2pt,
      align=center,
      text width=38mm,
      minimum height=14mm,
      inner sep=2mm,
      dashed
    },
    flow/.style={-{Latex[length=2.2mm]},semithick},
    fail/.style={-{Latex[length=2.2mm]},dashed},
    label/.style={font=\scriptsize,align=center,fill=white,inner sep=1pt}
  ]
    \node[stage] (select) at (0,0) {\textbf{Unity selection}\\
      stable entity/source ID\\and selection state};
    \node[gate] (identity) at (3.35,0) {\textbf{Gate 1: identity}\\
      pinned hash, unique asset, no ambiguity};
    \node[gate] (route) at (6.70,0) {\textbf{Gate 2: responsibility}\\
      asset $>$ group $>$ zone; local or direct};
    \node[gate] (binding) at (10.05,0) {\textbf{Gate 3: operation}\\
      exact scenario, provider, capability, binding, action};
    \node[stage] (model) at (13.40,0) {\textbf{Response model}\\
      called only after deterministic preflight};

    \draw[flow] (select) -- (identity);
    \draw[flow] (identity) -- (route);
    \draw[flow] (route) -- (binding);
    \draw[flow] (binding) -- (model);

    \node[reject] (reject) at (4.85,-3.25) {\textbf{Reject with location}\\
      stage + violated relation; no model call or retained mutation};
    \draw[fail] (identity.south) to[bend right=13]
      node[label,left,pos=.52] {stale/ambiguous} (reject.north west);
    \draw[fail] (route.south) --
      node[label,right,pos=.48] {no unique/permitted provider} (reject.north);
    \draw[fail] (binding.south) to[bend left=13]
      node[label,right,pos=.52] {zero/multiple bindings} (reject.north east);

    \node[gate] (evidence) at (13.40,-3.25) {\textbf{Gate 4: response}\\
      declared handoff and identifier agreement};
    \node[stage,text width=35mm] (present) at (9.30,-3.25)
      {\textbf{Commit and present}\\
      evidence retains target, provider, route, and action};
    \draw[flow] (model) -- (evidence);
    \draw[flow] (evidence) -- (present);

    \node[reject] (rollback) at (13.40,-5.65) {\textbf{Rollback}\\
      discard response-dependent proposal and provisional history};
    \draw[fail] (evidence) --
      node[label,right] {undeclared handoff or mismatch} (rollback);
  \end{tikzpicture}
  \caption{Fail-closed admission path.  Unity proposes a selected target, but
  the backend re-resolves identity, responsibility, and the exact operation
  from the pinned contract before generation.  Request-decidable failures stop
  before the response model; response-dependent failures roll back the
  proposal.}
  \Description{A two-row pipeline begins with a Unity selection and then passes
  through identity, responsibility, and operation gates before a response
  model.  Those preflight gates point to a rejection box.  The response then
  passes through an evidence gate before commit and presentation; a failed
  response points to a rollback box.}
  \label{fig:admission-boundary}
\end{figure*}

\subsection{Guarantee boundary}

The implementation provides a \emph{model-consistency guarantee}, not semantic
correctness.  If a deictic request is admitted, then under the pinned contract
it has one trusted asset, one uniquely resolved provider reachable in at most
one modeled handoff, and one exact capability-to-action binding; the accepted
response carries matching identifiers.  This statement does not establish that
the person intended the collider selected by the client, that the provider's
knowledge is correct, or that generated language answers the question well.
Those properties require trusted sensing, semantic ground truth, and human or
task-level evaluation outside the present contract.
